\section{Proof methods}
\label{sec:writing-foundations-04}

\subsection{Direct proof}

To prove \(A\to B\), assume \(A\) and derive \(B\).  In a metatheorem, the
assumption and conclusion must be stated at the same level.  Do not replace a
semantic hypothesis by a syntactic one unless an already proved theorem
licenses the replacement.

\subsection{Proof by contrapositive}

In classical metalogic, proving \(A\to B\) by proving
\(\neg B\to\neg A\) is legitimate.  In intuitionistic object logic, the
formula \((\neg B\to\neg A)\to(A\to B)\) is not generally derivable.  The
ambient logic must therefore be stated when this distinction matters.

\subsection{Proof by contradiction}

To prove \(A\) by contradiction in classical logic, assume \(\neg A\) and
derive \(\bottom\).  In intuitionistic logic this establishes
\(\neg\neg A\), and an additional principle is required to conclude \(A\).

\subsection{Induction}

State the induction domain, induction predicate, base cases, and induction
step.  Structural induction follows the grammar; induction on derivations
follows the final rule; induction on formula complexity follows the defined
complexity measure.  These methods may coincide extensionally in a particular
case, but the proof must identify the one used.

\subsection{Construction}

For an existence theorem, define the proposed witness and verify every required
property.  For a model-existence theorem, this includes a domain, an
interpretation for each symbol, and a satisfaction proof.  Naming a desired
object without constructing it or invoking an existence theorem is not a
construction.

\subsection{Counterexample}

To refute a universal statement, provide an object in the quantified domain and
verify that the predicate fails.  To refute semantic consequence for a
sentence argument, provide a countermodel.  For formulas with free variable
occurrences, provide a structure and variable assignment that satisfy the
premises and falsify the conclusion.  Either form of semantic refutation is
sufficient to refute derivability only after soundness of the stated calculus
has been proved.

